\documentclass[11pt]{article}
\usepackage[margin=1in]{geometry}
\usepackage{amsmath,amssymb,amsthm,hyperref}
\newtheorem{theorem}{Theorem}
\newtheorem{lemma}[theorem]{Lemma}
\newcommand{\MA}{\operatorname{MA}}
\newcommand{\R}{\mathbb R}
\title{Canonical generated energy and integrability for finite-dual measures}
\author{Research working paper, continuation 389}
\date{9 October 2026}
\begin{document}
\maketitle
\begin{abstract}
Under the exact bounded analytic premises and the newly proved, still
unreviewed bounded Green and tangent theorem of continuation 388, we
construct energy on pointwise decreasing limits of ordinary bounded
continuous potentials. Finite energy is independent of the decreasing
presentation, preserves order, and is concave with polynomial segment
energies. Every positive measure of finite full bounded-class dual energy
integrates every such generated finite-energy potential, with a uniform
energy bound and decreasing $L^1$ convergence. Every bounded analytic
Monge--Amp\`ere measure has finite full dual. These arguments concern
canonical generated energy and measure integrability. They do not define
unbounded Monge--Amp\`ere products or prove existence for arbitrary
finite-dual measures. All new proofs remain unreviewed.
\end{abstract}

\section{The exact bounded assumptions}
Use a compact Hausdorff analytic surface $X$ without boundary and the
entire ordinary class $\mathcal C_\theta$ of continuous bounded
$\theta$-psh potentials from continuation 388. Assume its ordinary
bounded positive mixed products, locality, raw mass $M>0$, and energy
normalization $E(0)=0$. The conclusions needed from the whole 388 proof
are convexity of $\mathcal C_\theta$, closure under adding real constants,
monotonicity and concavity of $E$, the constant rule $E(v+c)=E(v)+Mc$,
the complete-class tangent inequality, and its polynomial segment
first variation. These are explicitly unreviewed theorem dependencies,
not incoming accepted grants. In particular, we do not assume an
unbounded first-variation theorem.

For $u,v\in\mathcal C_\theta$ the bounded segment calculation also
gives the cocycle
\begin{equation}\label{eq:cocycle}
 E(u)-E(v)=\frac13\int_X(u-v)
       \bigl(\MA(u)+\theta_u\wedge\theta_v+\MA(v)\bigr).
\end{equation}
Indeed integrate $\int (u-v)\theta_{v+t(u-v)}^2$ from zero to one:
the three nonnegative quadratic Bernstein coefficients have integrals
$1/3$, including the coefficient $2t(1-t)$ of the mixed product.
Thus if $u\geq v$, each of the three weighted nonnegative integrals
is at most $3(E(u)-E(v))$. This is an identity of genuine bounded
products, not a definition of unbounded products.

\section{Canonical energy of decreasing presentations}
Call $u:X\to\R\cup\{-\infty\}$ generated finite-energy if
\[
                 u_j\in\mathcal C_\theta,\quad u_j\downarrow u,
                   \quad \inf_j E(u_j)>-\infty.
\]
Denote this class by $\mathcal E^1_{\downarrow}$.
It includes every bounded member by a constant sequence. Such a limit
is upper semicontinuous and locally $\theta$-psh in the smallest sheaf
closed under decreasing limits, because locally $\rho+u_j$ is psh and
decreases to $\rho+u$. No recognition of every function in the full
transfinite unbounded psh sheaf is asserted.

\begin{lemma}\label{lem:strict}
If continuous $v_k\downarrow v$ on a compact space and $a$ is
continuous with $v<a$ everywhere, then $v_k<a$ everywhere for some $k$.
The assertion includes $v(x)=-\infty$.
\end{lemma}
\begin{proof}
The open sets $\{v_k<a\}$ increase and cover $X$. A finite subcover
exists, and its largest index works on the whole space. No uniform
convergence to $v$ is claimed.
\end{proof}

\begin{theorem}\label{thm:canonical}
For $u\in\mathcal E^1_{\downarrow}$ the number
\[
                         E(u)=\lim_j E(u_j)
\]
is finite and independent of its decreasing bounded continuous
presentation. Every decreasing bounded continuous presentation of the
same $u$ has this same finite energy limit. The function $u$ is not
identically minus infinity. On the generated class, energy preserves
pointwise order and satisfies $E(u+c)=E(u)+Mc$.
\end{theorem}
\begin{proof}
Monotonicity makes $E(u_j)$ decrease; its lower bound gives a finite
limit. Let $v_k\downarrow u$ be any other bounded continuous
presentation, without assuming its energies bounded below. Fix $j$
and $\varepsilon>0$. Since $u<u_j+\varepsilon$ everywhere,
Lemma~\ref{lem:strict} gives $v_k<u_j+\varepsilon$ for large $k$.
Thus its energy limit, a priori possibly bottom, is at most
$E(u_j)+M\varepsilon$. Sending $j$ to infinity and then
$\varepsilon$ to zero gives $\lim E(v_k)\leq\lim E(u_j)$.
For each fixed $k$ apply the lemma instead to the sequence $u_j$
and the strict majorant $v_k+\varepsilon$.
Then $\lim E(u_j)\leq E(v_k)+M\varepsilon$; taking the infimum
over $k$ proves the reverse inequality and finiteness. This establishes
presentation independence without a diagonal flattening theorem.

If $u$ were bottom everywhere, the open sets $\{u_j<-A\}$ would
cover $X$ for every $A>0$. Compactness would give $u_j<-A$
everywhere for some $j$, hence $E(u_j)\leq-MA$. This contradicts
the fixed finite lower bound as $A$ increases.

For order preservation, suppose $u\leq v$ are generated. Fix a
bounded continuous term $v_k$ in a presentation of $v$. The strict
majorant $v_k+\varepsilon$ of $u$ and the lemma give
$E(u)\leq E(v_k)+M\varepsilon$. Take the infimum in $k$ and then
let $\varepsilon\downarrow0$. The constant rule follows by shifting
any presentation and taking its energy limit.
\end{proof}

\begin{theorem}\label{thm:segments}
The generated class is convex. Energy is concave on it. For any two
members $u,v$, the function $t\mapsto E((1-t)u+tv)$ on $[0,1]$ is
the restriction of a polynomial of degree at most three.
\end{theorem}
\begin{proof}
Choose presentations $u_j\downarrow u$, $v_j\downarrow v$ and
use $w_{j,t}=(1-t)u_j+tv_j$. It is a decreasing bounded continuous
presentation of the segment point. Bounded concavity gives
\[
 E(w_{j,t})\geq(1-t)E(u_j)+tE(v_j),
\]
uniformly bounded below for every fixed $t$. Thus the segment is
generated finite-energy and its canonical limit energy satisfies
the same concavity inequality. There is no ambiguous sum of opposite
infinities: both functions have only finite or negative-infinite values,
and the endpoints are understood separately.

Each $p_j(t)=E(w_{j,t})$ has degree at most three. At four distinct
points of $[0,1]$ its values converge to the finite canonical energies
just established. The invertible four-by-four Vandermonde system
therefore makes all coefficients of $p_j$ converge. The limiting
polynomial agrees with the canonical energy at every $t$, by
Theorem~\ref{thm:canonical}. Consequently the segment energies are
smooth as numerical functions of $t$. This argument does not identify
their derivatives with any unbounded analytic Monge--Amp\`ere operator.
\end{proof}

\section{Finite dual energy implies generated integrability}
For a positive finite Radon measure $\mu$ on $X$ put
\[
 E^*(\mu)=\sup_{v\in\mathcal C_\theta}
                      \left(E(v)-\int_Xv\,\mu\right).
\]
This is the full ordinary bounded-class dual. If it is finite then
$\mu(X)=M$: testing all constants gives
$E(c)-\int c\,\mu=c(M-\mu(X))$, unbounded above unless the
two masses agree. Also $E^*(\mu)\geq0$, by the zero test.

\begin{theorem}\label{thm:integrability}
If $E^*(\mu)<\infty$, every $u\in\mathcal E^1_{\downarrow}$ is
finite $\mu$-almost everywhere and belongs to $L^1(\mu)$. For
$u\leq0$ its precise estimate is
\begin{equation}\label{eq:bound}
                         \int_X(-u)\,\mu\leq E^*(\mu)-E(u).
\end{equation}
For every bounded continuous decreasing presentation $u_j\downarrow u$,
\[
                   \int_X|u_j-u|\,\mu\longrightarrow0.
\]
The right side of \eqref{eq:bound} gives a uniform negative-part
integrability bound on every nonpositive generated energy sublevel
$E(u)\geq-C$.
\end{theorem}
\begin{proof}
First subtract a constant from the whole presentation so all its terms
are nonpositive: its first term is bounded on compact $X$. For this
shifted presentation the definition of the bounded dual gives
\[
                    \int_X(-u_j)\,\mu\leq E^*(\mu)-E(u_j).
\]
The nonnegative functions $-u_j$ increase to $-u$. Monotone
convergence and canonical energy continuity prove the finite bound
\eqref{eq:bound}. In particular the pole set has zero $\mu$-mass.
Since $u$ is bounded above by its first continuous term, finiteness
of its negative-part integral also proves absolute integrability.

After this conclusion, $u_j-u\geq0$ and
\[
 \int_X|u_j-u|\,\mu=\int_Xu_j\,\mu-\int_Xu\,\mu\longrightarrow0
\]
by the same monotone convergence. The shift cancels in this difference.
For the stated estimate when the limit $u\leq0$ but its chosen
presentation need not be nonpositive, subtract one constant $B$ with
$u_j\leq B$, prove the estimate for $u-B$, and use
$\mu(X)=M$ and $E(u-B)=E(u)-MB$ to cancel $B$. Thus
\eqref{eq:bound} holds with the asserted normalization independently
of the presentation. Its uniform sublevel assertion follows immediately.
\end{proof}

\begin{theorem}\label{thm:bounded-measures}
For every $v\in\mathcal C_\theta$ the genuine bounded measure
$\MA(v)$ has finite full bounded dual and
\[
                      E^*(\MA(v))=E(v)-\int_Xv\,\MA(v).
\]
Consequently every generated finite-energy $u$ is integrable against
every bounded analytic $\MA(v)$, and every decreasing bounded
continuous presentation converges in the corresponding $L^1$ norm.
\end{theorem}
\begin{proof}
The entire-class tangent inequality at $v$ gives, for every bounded
test $w$,
\[
 E(w)-\int_Xw\,\MA(v)\leq E(v)-\int_Xv\,\MA(v).
\]
The right side is finite, and testing $w=v$ attains it. Now apply
Theorem~\ref{thm:integrability}. The measure is the actual ordinary
bounded product; no inverse or proposed unbounded operator is used.
\end{proof}

In particular, if continuation 388's geometric calibration and full
radial-dual theorem are independently verified, each of its specified
radial measures with finite cubic-tail integral integrates \emph{every}
generated finite-energy potential, including functions outside the
radial subclass. This follows from its full-class dual value and
Theorem~\ref{thm:integrability}; a restricted radial supremum alone
would not justify this assertion.

\section{Proof boundary and the remaining original problem}
All proofs above are new complete arguments under exactly the disclosed
bounded hypotheses and the whole unreviewed 388 theorem chain.
No additional primary body or recursive source proof is newly claimed
read here. The exact full primary readings, positivity scope, bounded
products, source-version qualifications and geometric calibration
dependencies are the entire 388 paper and its bound source-scope record,
which were read in full. No active native mathematical output has
been read or fed back, and no supplied/native program was executed
or imported. Full primary bodies, source pixels and whole book contexts
remain LOCALONLY.

We prove independence from bounded continuous decreasing presentations
and integrability for a canonical generated class. We do not assert
that it recognizes all intrinsic unbounded $\theta$-psh functions, that
every decreasing sequence of unbounded generated members has a bounded
continuous presentation obtained by an automatic diagonal argument,
or that the $L^1$ tests distinguish every pair of canonical potentials.
No separated compactness theorem is deduced from these integrability
estimates. Polynomial segment energies alone do not identify an analytic
Monge--Amp\`ere measure. Genuine unbounded mixed products, decreasing
continuity of those products, all-direction comparison, every obstacle
and orthogonality step, and existence/uniqueness for arbitrary finite-dual
measures remain to be proved for the full original target.

This paper is UNREVIEWED and does not close original P16.7 or increase
the original solved count. Counts remain 16 solved, 26 open, and 97
previously accepted scoped papers. Compilation verification is pending.
The protected P16.4 editor stays open and unchanged.
\end{document}
